\documentclass[11pt]{amsart}
\usepackage{calc,amssymb,amsmath,ulem,stmaryrd,fullpage}
\usepackage{enumerate}
\usepackage{alltt}
\RequirePackage[dvipsnames,usenames]{color}
\let\mffont=\sf
\normalem
%\input{kmacros3.sty}
\input{xy}
\xyoption{all}
\usepackage{hyperref}
\usepackage{graphicx}
\usepackage[all,cmtip]{xy}
\usepackage{verbatim}
\usepackage[left=0.95in,top=0.95in,right=0.95in,bottom=0.95in]{geometry}
\newcommand{\tauCohomology}{T}
\newcommand{\FCohomology}{S}
\usepackage{mathrsfs}


\theoremstyle{theorem}
%\newtheorem*{mainthm*}{Main Theorem}

\newcommand{\note}[1]{\marginpar{\sffamily\tiny #1}}

\def\todo#1{\textcolor{Mahogany}%
{\sffamily\footnotesize\newline{\color{Mahogany}\fbox{\parbox{\textwidth-15pt}{#1}}}\newline}}

\renewcommand{\O}{\mathcal O}
\newcommand{\ie}{{\itshape i.e.} }
\newcommand{\roundup}[1]{\lceil #1 \rceil}
\newcommand{\rounddown}[1]{\lfloor #1 \rfloor}
\renewcommand{\to}[1][]{\xrightarrow{\ #1\ }}
\newcommand{\onto}[1][]{\protect{\xrightarrow{\ #1\ }\hspace{-0.8em}\rightarrow}}
\newcommand{\into}[1][]{\lhook \joinrel \xrightarrow{\ #1\ }}
%\newcommand{DBDual}[1]{{\underline{\omega}_{#1}}}
\newcommand{\DBDual}[1]{{{\uuline \omega}{}^{\mydot}_{#1}}}
\newcommand{\Tt}{{\mathfrak{T}}}
%\newcommand{\bn}{{\mathfrak{n}} }
\newcommand{\sol}[1]{ \vskip 6pt{\color{blue} {\bf Solution:} #1} \vskip 6pt}

\newcommand{\bR}{{\mathbb{R}}}
\newcommand{\va}{{\vec{a}}}
\newcommand{\vb}{{\vec{b}}}
\newcommand{\vzero}{{\vec{0}}}
\newcommand{\vi}{{\vec{i}}}
\newcommand{\vj}{{\vec{j}}}
\newcommand{\vk}{{\vec{k}}}
\newcommand{\vh}{{\vec{h}}}
\newcommand{\vr}{{\vec{r}}}
\newcommand{\vu}{{\vec{u}}}
\newcommand{\vv}{{\vec{v}}}
\newcommand{\vw}{{\vec{w}}}
\newcommand{\vx}{{\vec{x}}}
\newcommand{\vy}{{\vec{y}}}
\newcommand{\vz}{{\vec{z}}}
\newcommand{\vT}{{\vec{T}}}
\newcommand{\vN}{{\vec{N}}}
\newcommand{\vF}{{\vec{F}}}
\newcommand{\vG}{{\vec{G}}}
\newcommand{\bF}{\mathbb{F}}
\newcommand{\bQ}{\mathbb{Q}}
\newcommand{\bZ}{\mathbb{Z}}
\newcommand{\bP}{\mathbb{P}}
\newcommand{\bA}{\mathbb{A}}
%\newcommand{\cF}{\mathcal{F}}
%\newcommand{\cG}{\mathcal{G}}
%\newcommand{\cH}{\mathcal{H}}
\newcommand{\image}{\mathrm{im}\,}
\newcommand{\coker}{\mathrm{coker}\,}
\newcommand{\Hom}{\mathrm{Hom}\,}
\newcommand{\depth}{\mathrm{depth}\,}
\newcommand{\Supp}{\mathrm{Supp}\,}
\newcommand{\Spec}{\mathrm{Spec}\,}
\newcommand{\Ann}{\mathrm{Ann}\,}
\newcommand{\sH}{\mathscr{H}\,}
\newcommand{\cI}{\mathcal{I}}
\newcommand{\cJ}{\mathcal{J}}
\newcommand{\cE}{\mathcal{E}}
\newcommand{\cF}{\mathcal{F}}
\newcommand{\cG}{\mathcal{G}}
\newcommand{\cH}{\mathcal{H}}
\newcommand{\cO}{\mathcal{O}}
\newcommand{\cM}{\mathcal{M}}
\newcommand{\cN}{\mathcal{N}}
\newcommand{\cK}{\mathcal{K}}
\newcommand{\fram}{\mathfrak{m}}
\newcommand{\Div}{\mathrm{Div}}

\begin{document}
\title{Worksheet \#5 -- Math 6140\\Spring 2019}
\author{Due Wednesday, April 3rd}

\maketitle
You may work in groups of up to 3.  Only one worksheet needs to be turned in per group.
\vskip 9pt
\noindent
The following exercises are essentially Hartshorne, III, Exercise 2.3.  \vskip 9pt
\noindent
Suppose that $X$ is a topological space and $Z \subseteq X$ is closed.  Consider the functor $\Gamma_Z(X, \bullet)$ which is defined by $\Gamma_Z(X, \cF) = \{s \in \Gamma(X, \cF)\;|\; \Supp(s) \subseteq Y \}$.\footnote{Recall that $\Supp(s) = \{x \in X\;|\;0\neq s_x \in \cF_x\}$.}  It is straightforward to see that this is a left exact functor.  The right derived functors are denoted by $H^i_Z(X, \bullet)$, and by the usual arguments, given a short exact sequence of Abelian groups,
\[
0 \to \cE \to \cF \to \cG \to 0
\]
we obtain a long exact sequence
\[
0 \to H^0_Z(X, \cE) \to H^0_Z(X, \cF) \to H^0_Z(X, \cG) \to H^1_Z(X, \cE) \to H^1_Z(X, \cF) \to H^1_Z(X, \cG) \to \dots
\]
\noindent
{\bf 1.}  Show that if $\cE$ is flasque, then $H^i_Z(X, \cE) = 0$ for $i > 0$.
\vskip 3pt
\emph{Hint:} Show that if $\cE$ is flasque in the above short exact sequence, then $H^0_Z(X, \cF) \twoheadrightarrow H^0_Z(X, \cG)$ surjects.
\newpage
\noindent
{\bf 2.}  Let $U = X \setminus Z$.  Show that for any $\cF$, there is a long exact sequence:
\[
0 \to H^0_Z(X, \cF) \to H^0(X, \cF) \to H^0(U, \cF) \to H^1_Z(X, \cF) \to H^1(X, \cF) \to H^1(U, \cF) \to \dots
\]
\vskip 3pt
\emph{Hint:} See worksheet 1, exercise \#8, and show that sequence from that exercise is exact if $\cF$ is flasque.  Also feel free to use anything from Chapter III, Section 1 of Hartshorne (or other sources for formal nonsense).
\vskip 250pt
\noindent
{\bf 3.}  Suppose that $A$ is a ring and that $I \subseteq A$ is a finitely generated ideal.  Consider the functor from $A$-modules to $A$-modules defined by $\Gamma_I(M) = \{m \in M \;|\; I^n m = 0 \text{ for some $n > 0$.} \}$.  Show that this functor is left exact.  Its right derived functors are denoted $H^i_I(\bullet)$ and are called local cohomology.
\vskip 200pt
\newpage
\noindent
{\bf 4.}  Suppose that $A$ is a Noetherian ring, $I \subseteq A$ is an ideal, and $M$ is an $A$-module.  Let $X = \mathrm{Spec}\; A$ and $Z = V(I) \subseteq X$.  Prove that
\[
H^i_I(M) = H^i_Z(X, \widetilde{M})
\]
for all $i > 0$.
\vskip 200pt
What follows are essentially exercises from Hartshorne chapter III section 3.
\vskip 3pt
Suppose that $A$ is a ring, $M$ is a finitely generated $A$-module and $I \subseteq A$ is an ideal.  Further suppose that $M \neq I M$.  Recall that the \emph{$I$-depth of $M$, or the depth\footnote{Some CA sources call this grade, Hartshorne calls it depth, as does the stacks project.  I'll stick with their terminology.  In the case of a local ring with $I = \fram$, these notions agree.} of $M$ along $I$} is the maximum length of a sequence of elements $\{x_1, \dots, x_d\} \subseteq I$ so that $x_{i+1}$ is not a zero divisor on $M/(x_1, \dots, x_{i})M$.  This is denoted by $\depth_I M$.
\vskip 9pt
\noindent
{\bf 5.}  For $A, I, M$ as above, show that $\depth_I M \geq k$ if and only if $H^i_I(M) = 0$ for all $i < k$.  
\vskip 3pt
\emph{Hint:} Use induction, the base case $k = 1$ is straightforward.
\newpage
Suppose $(X, \O_X)$ is a scheme.  A coherent sheaf $\cF$ of $\O_X$-modules is said to satisfy \emph{Serre's $\mathrm{S}_k$ property} if for each point $p \in X$ (including non-closed points), we have that $\depth_{\fram_p} \cF_p \geq \min(\dim \O_{X,p}, k)$.  Likewise, a finitely generated module $M$ over a ring $A$ is $\mathrm{S}_k$ if for each $Q \in \Spec A$, we have $\depth_Q M_Q \geq \min(\dim A_Q, k)$, this is equivalent to $\widetilde M$ satisfying $\mathrm{S}_k$ in the above sense.  \footnote{Warning, in some texts, $\dim \O_{X,p}$ or $\dim A_Q$ is replaced by $\dim \cF_p$ or $\dim M_Q$.}
\vskip 9pt
\noindent
{\bf 6.}  Suppose $X$ is a scheme and $Z \subseteq X$ is a closed subscheme of codimension\footnote{Codimension $\geq k$ means for every irreducible component $Z_i$ of $Z$, there exists a chain of irreducible subsets $Z_i = Y_0 \subseteq Y_1 \subseteq \dots Y_k \subseteq X$.} $\geq k$.  Show that if $\cF$ is $\mathrm{S}_k$, then $H^i_Z(X, \cF) = 0$ for all $i < k$.  If you just want to do the case where $k = 2$, that is fine too, and somewhat easier.  
\vskip 3pt
\emph{Hint:} Fix $k$ and use induction on $i$.  Fix a point $p \in Z$ (maybe even a generic point of $Z$) and consider the long exact sequence you get by from $H^i_Z(X, \bullet)$.  On the other hand, you can apply $H^i_p(X, \bullet)$ to this long exact sequence (or really to the short exact sequences you can break up the long exact sequence into).  This is a bit easier if you know spectral sequences, in which case just consider the spectral sequence obtained by composing the derived functors $H^j_p(X, \bullet)$ and $H^i_Z(X, \bullet)$.
\newpage
\noindent
{\bf 7.}  Conclude that if $\cF$ is $\mathrm{S}_2$ and $Z \subseteq X$ is codimension 2, then $\Gamma(X, \cF) = \Gamma(X \setminus Z, \cF)$.
\vskip 250pt
\noindent
{\bf 8.}  Suppose again that $X$ is a scheme, that $\cF$ is $\mathrm{S}_2$ and $Z \subseteq X$ is codimension 2.  Let $U = X \setminus Z$ and suppose that $i : U \to X$ is the (open) inclusion.  Show that $i_* \cF|_U \cong \cF$.  

\end{document}

